\section{Introduction}

Temporal knowledge graphs record facts as time-stamped quadruples $(s,r,o,t)$. They are useful for event streams, financial records, and other evolving collections in which an update affects later reasoning. The problem is that incoming facts do not fail in one uniform way. A fact may place an entity in an implausible relational role, preserve a plausible triple while moving it to an implausible time, or reveal that an expected fact is absent. A practical detector therefore needs more than a ranking score. It needs an explanation that identifies the historical evidence behind the score and a protocol that prevents unsupported rules from silently entering the detector.

AnoT addresses this setting with a detector--updater--monitor architecture that summarizes a \tkg into a rule graph and traverses that graph to score new facts \cite{Zhang_2024}. This is a useful alternative to opaque temporal embeddings: each prediction can be related to typed static rules and temporal rule edges. Yet the normal rule graph primarily answers whether a candidate resembles historical regularities. That question is insufficient when the candidate itself is locally plausible. A political statement, for example, may be a common event, while the same statement becomes suspicious because it is separated from its historical context by an atypical interval. Likewise, an unseen entity pair is weak evidence of a conceptual error if the entities have previously played compatible roles in the relation.

This paper extends the normal-rule perspective with error-aware evidence. We retain the original AnoT graph as the common backbone and learn a small, auditable evidence layer for each error type. For time errors, the layer is mined by comparing normal training contexts against temporally displaced training facts constructed under the original AnoT perturbation protocol. For conceptual errors, the layer measures whether an entity has historically appeared in the subject or object role of the observed relation. For missing errors, it retrieves train-only contextual patterns that support the expected target fact. These signals have different semantics and are therefore not forced into a single conflict definition.

We call the resulting framework \method (Conflict-Aware Rule Graph Orchestration for AnoT). The word ``agentic'' has a narrow meaning here. It refers to a typed, artifact-based workflow with explicit Candidate, Evidence, Validation, Scoring, Explanation, Evaluation, and Reflection components. It does not mean that a language model decides whether a fact is true. The current experiments use deterministic candidate sources and programmatic checkers; an optional LLM-candidate schema is not used to produce the reported metrics.

Our contributions are as follows.
\begin{itemize}[leftmargin=*,itemsep=2pt,topsep=2pt]
  \item We formulate an error-aware rule-graph interface that preserves AnoT's normal evidence while exposing distinct conflict and support channels for conceptual, time, and missing errors.
  \item We introduce contrastive temporal-gap mining for time errors. The miner creates training-only time-displaced references, selects atomic-rule gap patterns with MDL, enrichment, and Bayesian direction tests, and freezes the resulting rules before test evaluation.
  \item We implement an auditable rule-evidence workflow with typed artifacts, input/output hashes, deterministic validators, and evidence-grounded explanations. This workflow separates rule discovery from final scoring and makes the provenance of accepted rules inspectable.
\end{itemize}

\section{Background and problem setting}

Let $\mathcal{G}=\{(s,r,o,t)\}$ be a \tkg, where $s,o\in\mathcal{E}$, $r\in\mathcal{R}$, and $t$ is a discrete timestamp. In online anomaly detection, historical facts are used to build a detector, after which an arriving fact $x=(s,r,o,t)$ receives an anomaly score. A larger score indicates weaker compatibility with the learned graph. The rule-based formulation complements temporal knowledge graph forecasting methods that use temporal logical patterns for explainable prediction \cite{Liu_2022} and sits within the broader literature on temporal knowledge graph reasoning \cite{Cai_2023}.

We follow AnoT's three error categories \cite{Zhang_2024}. A conceptual error makes the entity--relation combination implausible. A time error preserves a plausible entity--relation combination but places it at an implausible timestamp. A missing error is a held-out true fact whose surrounding history should support its presence. The categories matter because their evidence directions differ: conceptual and time errors need evidence that raises suspicion, whereas a missing target should receive a higher score when contextual support is present.

AnoT maps a fact to typed static rules of the form $(\tau_s,r,\tau_o,d)$, where $\tau_s$ and $\tau_o$ are entity types and $d\in\{\mathrm{in},\mathrm{out}\}$ is an orientation. Temporal edges connect static rules that occur in historical sequence. The graph is learned with a minimum-description-length objective and supplies reachable rules as explanations during scoring \cite{Zhang_2024}. \method uses this graph unchanged as the normal-evidence backbone. This choice is deliberate: conventional rule mining seeks compact, high-support regularities \cite{Galarraga_2015}, whereas our additional layer asks a different question---which rule contexts distinguish a specific error mechanism from normal updates.

For an evaluation item $x$, let $y_k(x)\in\{0,1\}$ denote its label for error type $k\in\{\mathrm{conceptual},\mathrm{time},\mathrm{missing}\}$. The label direction follows the original benchmark: a corrupted conceptual or temporal update has label one, while a held-out true fact has label one for the missing-fact task. Thus, the same graph observation can have different semantics across tasks. A long gap is evidence against the plausibility of a time-stamped update, but contextual recurrence is evidence \emph{for} an omitted true fact. This distinction is why we retain three evidence channels rather than relabeling all graph signals as ``conflicts.''

